%! AP Statistics — Unit 1, Lesson 1.4 — Group Activity
\documentclass[10pt]{article}
\usepackage{apstats-article}
\usepackage{apstats-boxes}

%=============================================================================
\begin{document}

\pageheader{Unit 1, Lesson 1.4}{Group Activity: Bar Charts \& Misleading Graphs}
\namepartnerperiod

\vspace{0.1in}

\begin{tcolorbox}[colback=sky, colframe=navy, boxrule=0.8pt, arc=2mm,
    left=3mm, right=3mm, top=2mm, bottom=2mm]
\small
\textbf{Part 1} (group, 8 min): Complete the table and construct a relative
frequency bar chart. Then write an interpretation.\\[4pt]
\textbf{Part 2} (group, 5 min): Identify and describe a misleading version.
Optional: compute sector angles for a pie chart.\\[4pt]
\textbf{Part 3} (individual, 3 min): Reflection question.
\end{tcolorbox}

\vspace{0.14in}

% ── PART 1 ───────────────────────────────────────────────────────────────────
{\large\bfseries\color{navy} Part 1 \quad Construct the Bar Chart}

\vspace{0.1in}

\noindent\small A survey of \textbf{50 AP Statistics students} asks: ``Which streaming
service do you use most?'' Results are below.

\vspace{0.1in}

\renewcommand{\arraystretch}{2.0}
\begin{tabularx}{\linewidth}{@{} >{\bfseries}p{0.22\linewidth}
                                    C{0.18\linewidth}
                                    C{0.22\linewidth}
                                    C{0.22\linewidth} @{}}
\rowcolor{sky}
\textbf{Service} & \textbf{Frequency} & \textbf{Rel.\ Frequency} & \textbf{Angle (${}^\circ$)} \\
\midrule
Netflix  & 21 & \blank{2.5cm} & \blank{2cm} \\
Disney+  & 10 & \blank{2.5cm} & \blank{2cm} \\
Hulu     & 9  & \blank{2.5cm} & \blank{2cm} \\
HBO Max  & 6  & \blank{2.5cm} & \blank{2cm} \\
Other    & 4  & \blank{2.5cm} & \blank{2cm} \\
\midrule
\rowcolor{sky}\textbf{Total} & \blank{2cm} & \blank{2.5cm} & \blank{2cm} \\
\end{tabularx}

\vspace{0.14in}

\noindent\textbf{Construct a relative frequency bar chart below.}\\[2pt]
\small Include: (1) a title, (2) a labeled x-axis with categories, (3) a labeled
y-axis starting at 0, (4) gaps between bars.

\vspace{8pt}
\begin{tcolorbox}[colback=white, colframe=navy, boxrule=0.8pt, arc=2mm,
    left=3mm, right=3mm, top=2mm, bottom=2mm, height=2.8in]
\end{tcolorbox}

\vspace{0.12in}

\begin{enumerate}[leftmargin=1.5em, itemsep=8pt]
  \item Does your relative frequency column sum to 1? \blank{1cm}
        If not, find and fix the rounding error.

  \item Write a complete sentence identifying the most popular and least popular
        streaming service, including relative frequencies and context.\\[5pt]
        \writeline\\[1pt]\writeline

  \item Write a second sentence describing the overall pattern of the distribution
        (e.g., is one service much more popular than the others?).\\[5pt]
        \writeline\\[1pt]\writeline
\end{enumerate}

\vspace{0.14in}

% ── PART 2 ───────────────────────────────────────────────────────────────────
{\large\bfseries\color{navy} Part 2 \quad Misleading Graphs \& Pie Chart (Optional)}

\vspace{0.1in}

\begin{enumerate}[leftmargin=1.5em, itemsep=8pt]

  \item Suppose someone redraws your bar chart with the y-axis starting at 30\%
        instead of 0\%. Describe how this changes what a reader sees. Which service
        would appear most dramatically more popular?\\[5pt]
        \writeline\\[1pt]\writeline\\[1pt]\writeline

  \item A student presents a 3-D pie chart of this data with the Netflix slice
        tilted toward the viewer. Explain why this display might be misleading
        even though the data are correct.\\[5pt]
        \writeline\\[1pt]\writeline

  \item \textbf{Optional --- Pie Chart:} Using the angles you computed in the table,
        describe the largest and smallest sectors. Is a bar chart or pie chart more
        effective for this dataset? Why?\\[5pt]
        \writeline\\[1pt]\writeline

\end{enumerate}

\vspace{0.14in}

% ── PART 3 ───────────────────────────────────────────────────────────────────
{\large\bfseries\color{navy} Part 3 \quad Individual Reflection}

\vspace{0.1in}

\begin{tcolorbox}[colback=goldbg, colframe=goldacc, boxrule=0.8pt, arc=2mm,
    left=3mm, right=3mm, top=2mm, bottom=2mm]
\small\textbf{In your own words:} What is the single most important thing to check
when you look at a bar chart made by someone else? Why does it matter?
Give a specific example from today's activity.\\[6pt]
\writeline\\[1pt]\writeline\\[1pt]\writeline
\end{tcolorbox}

\end{document}
